\documentclass[12pt]{article}
\usepackage[utf8]{inputenc}
\usepackage[T1]{fontenc}
\usepackage[margin=1in]{geometry}
\usepackage{setspace}
\usepackage{xcolor}
\usepackage[colorlinks=true, urlcolor=blue!60!black, linkcolor=blue!60!black]{hyperref}
\setlength{\parindent}{2em}
\setlength{\parskip}{1.25em}
\renewcommand{\baselinestretch}{1.0}
\begin{document}
\begin{center}
{\LARGE \textit{Humphrey's Executor v. United States}, 295 U.S. 602 (1935)\par}
\vspace{0.8em}
{\large Case Brief\par}
\end{center}
\vspace{1em}

\section{Memory Jogger}

In a landmark 1935 separation-of-powers decision (later overruled in 2026 by \textit{Trump v. Slaughter}), a unanimous Supreme Court held that Congress could constitutionally restrict the President to ``for-cause'' removal of Federal Trade Commissioners because the FTC was an independent, nonpartisan body performing ``quasi-legislative'' and ``quasi-judicial'' duties rather than ``purely executive'' functions.

\section{Detailed Case Facts}

On December 10, 1931, President Herbert Hoover nominated William E.\ Humphrey, a conservative Republican, to succeed himself as a member of the Federal Trade Commission (FTC); the Senate confirmed him, and he was commissioned for a seven-year term set to expire on September 25, 1938 (\href{https://supreme.justia.com/cases/federal/us/295/602/}{295 U.S.\ at 618}; \href{https://www.courtlistener.com/opinion/102475/humphreys-v-united-states/}{CourtListener}). Section 1 of the Federal Trade Commission Act of 1914 (38 Stat.\ 717, 15 U.S.C.\ \S\S\ 41--42) created a five-member commission with staggered seven-year terms, mandated that no more than three commissioners belong to the same political party, and provided that ``[a]ny commissioner may be removed by the President for inefficiency, neglect of duty, or malfeasance in office'' (\href{https://supreme.justia.com/cases/federal/us/295/602/}{295 U.S.\ at 619--20}). Under \S\S\ 5--7 of the Act, the FTC was empowered to prevent ``unfair methods of competition in commerce'' through complaint, notice, hearing, and cease-and-desist orders reviewable in federal circuit courts, to conduct investigations and submit reports to Congress, and to serve as a master in chancery in antitrust equity suits referred by federal courts (\href{https://supreme.justia.com/cases/federal/us/295/602/}{295 U.S.\ at 620--21}).

After taking office in 1933, President Franklin D.\ Roosevelt clashed with Humphrey over New Deal regulatory policies. On July 25, 1933, Roosevelt wrote to Humphrey requesting his resignation on the ground ``that the aims and purposes of the Administration with respect to the work of the Commission can be carried out most effectively with personnel of my own selection,'' while explicitly disclaiming any reflection on Humphrey personally or on his services (\href{https://supreme.justia.com/cases/federal/us/295/602/}{295 U.S.\ at 618}). After Humphrey requested time to consult his friends, Roosevelt wrote again on August 31, 1933, stating: ``You will, I know, realize that I do not feel that your mind and my mind go along together on either the policies or the administering of the Federal Trade Commission, and, frankly, I think it is best for the people of this country that I should have a full confidence'' (\href{https://supreme.justia.com/cases/federal/us/295/602/}{295 U.S.\ at 619}).

Humphrey refused to resign. On October 7, 1933, President Roosevelt sent him a one-sentence dismissal notice: ``Effective as of this date, you are hereby removed from the office of Commissioner of the Federal Trade Commission'' (\href{https://supreme.justia.com/cases/federal/us/295/602/}{295 U.S.\ at 619}). Roosevelt cited none of the statutory causes listed in \S\ 1 of the FTC Act. Humphrey never acquiesced in the firing and continued to insist he remained a commissioner entitled to his \$10,000 annual salary until his death on February 14, 1934 (\href{https://supreme.justia.com/cases/federal/us/295/602/}{295 U.S.\ at 618--19}).

\section{Procedural History}

Following Humphrey's death, Samuel F.\ Rathbun, the executor of Humphrey's estate (hence the official docket title \textit{Rathbun, Executor v. United States}), filed suit against the United States in the U.S.\ Court of Claims to recover Humphrey's unpaid salary from October 8, 1933, through February 14, 1934 (\href{https://supreme.justia.com/cases/federal/us/295/602/}{295 U.S.\ at 618, 632 n.*}; \href{https://www.oyez.org/cases/1900-1940/295us602}{Oyez}). Pursuant to \S\ 3(a) of the Act of February 13, 1925 (43 Stat.\ 936, 939; 28 U.S.C.\ \S\ 288), the Court of Claims certified two questions of law directly to the U.S.\ Supreme Court (\href{https://supreme.justia.com/cases/federal/us/295/602/}{295 U.S.\ at 618--19}):
\begin{enumerate}
    \item Did \S\ 1 of the FTC Act restrict the President's power to remove a commissioner solely to the enumerated causes (``inefficiency, neglect of duty, or malfeasance in office'')?
    \item If so, was that statutory restriction valid under the U.S.\ Constitution?
\end{enumerate}
The Supreme Court heard oral argument on May 1, 1935, and issued its unanimous decision answering both certified questions ``Yes'' on May 27, 1935---the same ``Black Monday'' on which the Court struck down the National Industrial Recovery Act in \textit{A.L.A. Schechter Poultry Corp. v. United States} (\href{https://supreme.justia.com/cases/federal/us/295/602/}{295 U.S.\ at 602, 632}; \href{https://www.columbialawreview.org/content/quasi-judicial-a-history-and-tradition/}{Columbia Law Review}).

\section{Judicial Votes}

The Supreme Court answered both certified questions in the affirmative by a vote of \textbf{9--0} (\href{https://supreme.justia.com/cases/federal/us/295/602/}{295 U.S.\ at 618, 632}; \href{https://caselaw.findlaw.com/court/us-supreme-court/295/602.html}{FindLaw}):
\begin{itemize}
    \item \textbf{Opinion of the Court (8):} Justice George Sutherland delivered the opinion of the Court, joined by Chief Justice Charles Evans Hughes and Justices Willis Van Devanter, Louis D.\ Brandeis, Pierce Butler, Harlan Fiske Stone, Owen J.\ Roberts, and Benjamin N.\ Cardozo.
    \item \textbf{Separate Concurring Statement (1):} Justice James Clark McReynolds agreed that both certified questions should be answered in the affirmative, appending a brief statement referencing his dissenting opinion in \textit{Myers v. United States}, 272 U.S.\ 52, 178 (1926), for his broader views on the President's removal power (\href{https://caselaw.findlaw.com/court/us-supreme-court/295/602.html}{295 U.S.\ at 632}).
    \item \textbf{Dissenting Opinions (0):} None.
\end{itemize}

\section{Holding}

\textbf{Judgment:} Certified Question No.\ 1 answered ``Yes''; Certified Question No.\ 2 answered ``Yes'' (\href{https://supreme.justia.com/cases/federal/us/295/602/}{295 U.S.\ at 632}). President Roosevelt lacked both statutory and constitutional authority to remove Commissioner Humphrey for policy disagreements without establishing one of the causes specified in \S\ 1 of the FTC Act. \textit{(Note: On June 29, 2026, the Supreme Court formally overruled Humphrey's Executor in a 6--3 decision in \href{https://supreme.justia.com/cases/federal/us/609/25-332/}{Trump v. Slaughter, 609 U.S.\ 422 (2026)}.)}

\textbf{Rule Announced:} First, as a matter of statutory interpretation, when Congress fixes a definite term of office for members of an independent, nonpartisan commission and enumerates specific causes for removal, Congress intends to prohibit removal during that term except for the enumerated causes (\href{https://supreme.justia.com/cases/federal/us/295/602/}{295 U.S.\ at 623--26}). Second, as a matter of constitutional law, whether the President possesses an unrestrictable power to remove an officer depends on ``the character of the office'': \textit{Myers v. United States}, 272 U.S.\ 52 (1926), guarantees exclusive, illimitable presidential removal power only over ``purely executive officers,'' whereas Congress may constitutionally fix terms and impose for-cause removal protections on officers of independent agencies performing ``quasi-legislative'' and ``quasi-judicial'' functions (\href{https://supreme.justia.com/cases/federal/us/295/602/}{295 U.S.\ at 627--32}).

\textbf{Holding vs.\ Dicta:}
\begin{itemize}
    \item \textbf{Binding Holding (as decided in 1935):} Section 1 of the FTC Act restricted removal of FTC commissioners to the three enumerated causes, and that statutory restriction was constitutional because FTC commissioners performed predominantly quasi-legislative and quasi-judicial duties rather than purely executive duties. The Court expressly limited \textit{Myers} to purely executive officers and disapproved broader dicta in \textit{Myers} suggesting an unlimited presidential removal power over all civil officers (\href{https://supreme.justia.com/cases/federal/us/295/602/}{295 U.S.\ at 626--28}).
    \item \textbf{Dicta / Reserved Questions:} Justice Sutherland expressly declined to resolve the ``field of doubt'' lying between purely executive officers (governed by \textit{Myers}) and officers of quasi-legislative/quasi-judicial agencies like the FTC, leaving intermediate offices ``for future consideration and determination as they may arise'' (\href{https://supreme.justia.com/cases/federal/us/295/602/}{295 U.S.\ at 632}). Similarly, the Court's observations regarding the Interstate Commerce Commission and the legislative Court of Claims were illustrative analogies rather than holdings on those specific bodies (\href{https://supreme.justia.com/cases/federal/us/295/602/}{295 U.S.\ at 629}).
\end{itemize}

\section{Analysis of Opinions}

\subsection{Majority Opinion (Justice Sutherland)}

Justice Sutherland divided the Court's opinion into two main parts corresponding to the two certified questions (\href{https://supreme.justia.com/cases/federal/us/295/602/}{295 U.S.\ at 621--32}; \href{https://tile.loc.gov/storage-services/service/ll/usrep/usrep295/usrep295602/usrep295602.pdf}{U.S.\ Reports PDF}):

\textbf{1.\ Statutory Construction (\S\ 1 of the FTC Act):} The Government argued under \textit{Shurtleff v. United States}, 189 U.S.\ 311 (1903), that listing removal causes did not exclude at-will presidential removal. Justice Sutherland distinguished \textit{Shurtleff}: the statute there fixed \textit{no} term of office for general customs appraisers, meaning that reading the causes as exclusive would have conferred unprecedented life tenure on a non-judicial civil officer (\href{https://supreme.justia.com/cases/federal/us/295/602/}{295 U.S.\ at 621--23}). By contrast, the FTC Act fixed staggered seven-year terms, and the Senate committee reports and floor debates---properly consulted to illuminate statutory purpose and the evils Congress sought to remedy---confirmed that Congress intended an expert, nonpartisan commission ``independent of executive authority, except in its selection'' (\href{https://supreme.justia.com/cases/federal/us/295/602/}{295 U.S.\ at 623--26}).

\textbf{2.\ Constitutional Validity (Separation of Powers):} Turning to Article II, Justice Sutherland narrowed \textit{Myers v. United States}, 272 U.S.\ 52 (1926), invoking Chief Justice Marshall's maxim in \textit{Cohens v. Virginia}, 19 U.S.\ (6 Wheat.) 264, 399 (1821), that general expressions beyond the point actually decided do not bind future courts under \textit{stare decisis} (\href{https://supreme.justia.com/cases/federal/us/295/602/}{295 U.S.\ at 626--27}). \textit{Myers} involved a first-class postmaster---a ``purely executive'' subordinate restricted to executive duties. The FTC, by contrast, was an administrative body created to carry out legislative policies, investigate and report to Congress (\S\ 6), adjudicate unfair methods of competition (\S\ 5), and act as a master in chancery for courts (\S\ 7) (\href{https://supreme.justia.com/cases/federal/us/295/602/}{295 U.S.\ at 627--28}). It could not ``in any proper sense be characterized as an arm or an eye of the executive''; any executive function it exercised was merely incidental to its quasi-legislative or quasi-judicial powers (\href{https://supreme.justia.com/cases/federal/us/295/602/}{295 U.S.\ at 628 \& n.*}).

Sutherland then grounded Congress's power to restrict removal in the separation of powers itself: because ``one who holds his office only during the pleasure of another cannot be depended upon to maintain an attitude of independence against the latter's will,'' giving the President illimitable removal power over legislative and judicial aids would allow the Executive to coerce the other branches---violating the maxim that making one ``master in his own house precludes him from imposing his control in the house of another who is master there'' (\href{https://supreme.justia.com/cases/federal/us/295/602/}{295 U.S.\ at 629--30}). Finally, Sutherland reread the ``Decision of 1789'' and \textit{Marbury v. Madison}, 5 U.S.\ (1 Cranch) 137 (1803), as distinguishing purely executive officers (like the Secretary of Foreign Affairs, whose ``acts are his acts'') from officers with judicial or non-executive duties (like the Comptroller of the Treasury or a District of Columbia justice of the peace) (\href{https://supreme.justia.com/cases/federal/us/295/602/}{295 U.S.\ at 630--31}).

\subsection{Concurring Opinions}

\textbf{Justice McReynolds (Separate Statement):} Justice McReynolds did not write a full concurrence, but appended a two-sentence statement at the end of the opinion agreeing that both certified questions should be answered ``Yes'' and referring to his dissenting opinion in \textit{Myers v. United States}, 272 U.S.\ 52, 178 (1926), for his broader view that the Constitution does not grant the President an inherent, illimitable power to remove Senate-confirmed statutory officers (\href{https://caselaw.findlaw.com/court/us-supreme-court/295/602.html}{295 U.S.\ at 632}).

\subsection{Dissenting Opinions}

None.

\section{Examples: Future Applications}

\textit{Note: The classifications below reflect how future disputes were analyzed under the Humphrey's Executor framework (as subsequently narrowed by \href{https://www.supremecourt.gov/opinions/19pdf/19-7_n6io.pdf}{Seila Law LLC v. CFPB, 591 U.S.\ 197 (2020)}), prior to the Supreme Court's June 29, 2026 decision overruling Humphrey's Executor in \href{https://supreme.justia.com/cases/federal/us/609/25-332/}{Trump v. Slaughter, 609 U.S.\ 422 (2026)}.}

\begin{enumerate}
    \item \textbf{Same-Side (For-Cause Removal Upheld under \textit{Humphrey's Executor}): Multi-Member, Bipartisan Adjudicatory Commission.} Suppose Congress creates a five-member War Claims Commission or a bipartisan multi-member regulatory tribunal with staggered terms tasked primarily with adjudicating statutory claims and issuing reports to Congress, without direct partisan policymaking. Under \textit{Humphrey's Executor} (and as applied in \textit{Wiener v. United States}, 357 U.S.\ 349 (1958)), such officers perform predominantly quasi-judicial duties requiring independence from executive coercion, so Congress may prohibit at-will presidential removal.

    \item \textbf{Same-Side (For-Cause Removal Upheld under \textit{Humphrey's Executor}): Multi-Member Expert Regulatory Board with Balanced Partisanship.} Suppose Congress establishes a five-member Securities or Labor Board with staggered fixed terms, a statutory partisan-balance requirement (no more than three members from one party), and removal limited to ``inefficiency, neglect of duty, or malfeasance in office.'' Under \textit{Humphrey's Executor}'s classic blueprint, the agency's multi-member expert structure and mix of rulemaking and adjudication fit squarely within the exception to \textit{Myers}---though after \href{https://supreme.justia.com/cases/federal/us/609/25-332/}{\textit{Trump v. Slaughter} (2026)}, multi-member agencies that execute congressional mandates against private parties (through rulemaking, adjudication, or enforcement suits) are now subject to at-will presidential removal (\href{https://ballotpedia.org/Humphrey\%27s_Executor_v._United_States}{Ballotpedia}).

    \item \textbf{Opposite-Side (For-Cause Removal Struck Down): Cabinet Secretary or Purely Executive Officer.} Suppose Congress enacts a statute prohibiting the President from removing the Secretary of State, the Secretary of Defense, or a regional Postmaster except for ``malfeasance in office.'' Under both \textit{Myers} and \textit{Humphrey's Executor}, such officials are ``purely executive officers'' acting as arms and eyes of the President; restricting their removal unconstitutionally interferes with the President's Article II Vesting and Take Care powers (\href{https://supreme.justia.com/cases/federal/us/295/602/}{295 U.S.\ at 627--28}).

    \item \textbf{Opposite-Side (For-Cause Removal Struck Down): Single-Director Independent Agency Wielding Substantial Executive Power.} Suppose Congress creates an agency headed by a single Director serving a five-year term, removable only for ``inefficiency, neglect of duty, or malfeasance in office,'' with unilateral power to promulgate binding regulations, bring federal enforcement actions, and seek civil penalties. As the Supreme Court held in \href{https://www.supremecourt.gov/opinions/19pdf/19-7_n6io.pdf}{\textit{Seila Law LLC v. CFPB}, 591 U.S.\ 197 (2020)}, such a structure falls outside \textit{Humphrey's Executor} because it vests significant executive power in a single unaccountable individual rather than a balanced multi-member body (\href{https://harvardlawreview.org/print/vol-137/the-presumption-against-novelty-in-the-roberts-courts-separation-of-powers-case-law/}{Harvard Law Review}).

    \item \textbf{Fence-Sitter (Close / Contested Case): Federal Reserve Board of Governors.} Suppose Congress insulates members of the Board of Governors of the Federal Reserve System with 14-year staggered terms and for-cause removal protection. Even as the Roberts Court dismantled \textit{Humphrey's Executor} for conventional regulatory commissions in \href{https://supreme.justia.com/cases/federal/us/609/25-332/}{\textit{Trump v. Slaughter} (2026)}, the Federal Reserve occupied a uniquely contested middle ground due to its historically distinct, quasi-private central-banking lineage and macroeconomic sensitivity---leading the Court on the very same day in \textit{Trump v. Cook} (2026) to carve out Federal Reserve independence from \textit{Slaughter}'s general at-will removal rule (\href{https://www.dlapiper.com/en/insights/publications/2026/07/supreme-court-overrules-humphreys-executor-upholds-federal-reserve-independence}{DLA Piper}; \href{https://en.wikipedia.org/wiki/Trump_v._Slaughter}{Wikipedia}).
\end{enumerate}

\section{Critique}

\textbf{Originalist and Textualist Perspectives:} For decades, unitary-executive theorists and originalist jurists mounted a sustained critique of \textit{Humphrey's Executor}. Critics---including Justice Scalia dissenting in \textit{Morrison v. Olson}, 487 U.S.\ 654, 724--27 (1988), Justice Thomas concurring in \href{https://www.supremecourt.gov/opinions/19pdf/19-7_n6io.pdf}{\textit{Seila Law}}, and scholars such as Steven Calabresi and Saikrishna Prakash---argued that Article II's Vesting Clause (``The executive Power shall be vested in a President'') and Take Care Clause vest \textit{all} executive power in a single, electorally accountable President, leaving no constitutional home for a headless ``fourth branch'' (\href{https://www.gwlr.org/wp-content/uploads/2016/01/83-Geo-Wash-L-Rev-1835.pdf}{George Washington Law Review, \textit{Debunking Humphrey's Executor}}; \href{https://www.cato.org/supreme-court-review/2019-2020/removal-power-critical-guide}{Cato Supreme Court Review}). Originalist critics also attacked Sutherland's ``quasi-legislative'' and ``quasi-judicial'' labels as ``handwaving and obfuscating phrases'' and contended that even in 1935---and certainly by the twenty-first century---the FTC exercised substantial executive law-enforcement power, meaning \textit{Humphrey's Executor} did not even satisfy its own exception (\href{https://www.supremecourt.gov/DocketPDF/25/25-332/384619/20251114165832949_25-332-Brief\%20of\%20Amici.pdf}{Amici Brief in \textit{Trump v. Slaughter}}; \href{https://ndlawreview.org/interring-the-unitary-executive/}{Notre Dame Law Review}). This critique ultimately prevailed when the Supreme Court overruled \textit{Humphrey's Executor} in \href{https://supreme.justia.com/cases/federal/us/609/25-332/}{\textit{Trump v. Slaughter}, 609 U.S.\ 422 (2026)}, holding that an agency executing a congressional mandate against private parties exercises executive power and must remain subject to at-will presidential removal (\href{https://www.theusconstitution.org/litigation/slaughter-v-trump/}{Constitutional Accountability Center}).

\textbf{Progressive, Historical, and Doctrinal Perspectives:} Progressive scholars, administrative-law historians, and a growing wave of originalist-method critics of the unitary executive---including Christine Kexel Chabot, Jed Shugerman, Beau Baumann, and Peter Shane---have defended \textit{Humphrey's Executor}'s core premise on both historical and functional grounds (\href{https://ndlawreview.org/interring-the-unitary-executive/}{Notre Dame Law Review}; \href{https://www.yalejreg.com/nc/quasi-judicial-a-history-and-tradition-by-beau-j-baumann-jed-h-shugerman/}{Yale Journal on Regulation}; \href{https://www.acslaw.org/seila-law-lcc-v-cfpb-and-the-persistent-myths-of-presidentialism/}{ACS}). Historical research demonstrates that the Founding generation and the First Congress did not embrace an absolute presidential removal power, having structured early financial and adjudicatory offices (such as the Sinking Fund Commission and the Comptroller of the Treasury) with significant independence, and that ``quasi-judicial'' independence reflects a long Anglo-American tradition designed to guarantee impartial adjudication and due process (\href{https://www.columbialawreview.org/content/quasi-judicial-a-history-and-tradition/}{Columbia Law Review}). Doctrinally and functionally, defenders argue that insulating expert, multi-member agencies from partisan firing protects regulatory stability, technical expertise, and checks and balances against executive aggrandizement---a point echoed by Justice Sotomayor's dissent in \textit{Trump v. Slaughter}, which warned that overturning \textit{Humphrey's Executor} undid 90 years of precedent and 140 years of bipartisan political practice (\href{https://www.theusconstitution.org/litigation/slaughter-v-trump/}{Constitutional Accountability Center}; \href{https://www.americanprogress.org/article/what-is-humphreys-executor-and-why-should-you-care-about-it/}{Center for American Progress}).

\textbf{Underlying Competing Values:} \textit{Humphrey's Executor} sits at the heart of the clash between \textbf{democratic accountability and executive unity}---the belief that all officials who wield federal law-enforcement power must answer directly to the elected President so voters know whom to blame---and \textbf{institutional independence, expertise, and anti-autocracy}---the belief that Congress must be able to insulate technical regulators and administrative adjudicators from short-term partisan pressure and presidential coercion.

\section{Key Quotations}

\begin{quote}
``Putting aside dicta, which may be followed if sufficiently persuasive but which are not controlling, the necessary reach of the decision goes far enough to include all purely executive officers. It goes no farther; much less does it include an officer who occupies no place in the executive department, and who exercises no part of the executive power vested by the Constitution in the President.''
\end{quote}
\noindent---Justice Sutherland, \href{https://supreme.justia.com/cases/federal/us/295/602/}{\textit{Humphrey's Executor v. United States}, 295 U.S.\ 602, 627--28 (1935)} (\href{https://www.courtlistener.com/opinion/102475/humphreys-v-united-states/}{CourtListener}).

\begin{quote}
``Such a body cannot in any proper sense be characterized as an arm or an eye of the executive. Its duties are performed without executive leave, and, in the contemplation of the statute, must be free from executive control.''
\end{quote}
\noindent---Justice Sutherland, \href{https://supreme.justia.com/cases/federal/us/295/602/}{\textit{Humphrey's Executor v. United States}, 295 U.S.\ 602, 628 (1935)} (\href{https://www.courtlistener.com/opinion/102475/humphreys-v-united-states/}{CourtListener}).

\begin{quote}
``The authority of Congress, in creating quasi-legislative or quasi-judicial agencies, to require them to act in discharge of their duties independently of executive control cannot well be doubted, and that authority includes, as an appropriate incident, power to fix the period during which they shall continue in office, and to forbid their removal except for cause in the meantime. For it is quite evident that one who holds his office only during the pleasure of another cannot be depended upon to maintain an attitude of independence against the latter's will.''
\end{quote}
\noindent---Justice Sutherland, \href{https://supreme.justia.com/cases/federal/us/295/602/}{\textit{Humphrey's Executor v. United States}, 295 U.S.\ 602, 629 (1935)} (\href{https://www.courtlistener.com/opinion/102475/humphreys-v-united-states/}{CourtListener}).

\begin{quote}
``The fundamental necessity of maintaining each of the three general departments of government entirely free from the control or coercive influence, direct or indirect, of either of the others has often been stressed, and is hardly open to serious question. \ldots\ The sound application of a principle that makes one master in his own house precludes him from imposing his control in the house of another who is master there.''
\end{quote}
\noindent---Justice Sutherland, \href{https://supreme.justia.com/cases/federal/us/295/602/}{\textit{Humphrey's Executor v. United States}, 295 U.S.\ 602, 629--30 (1935)} (\href{https://www.courtlistener.com/opinion/102475/humphreys-v-united-states/}{CourtListener}).

\begin{quote}
``Whether the power of the President to remove an officer shall prevail over the authority of Congress to condition the power by fixing a definite term and precluding a removal except for cause will depend upon the character of the office; the \textit{Myers} decision, affirming the power of the President alone to make the removal, is confined to purely executive officers, and, as to officers of the kind here under consideration, we hold that no removal can be made during the prescribed term for which the officer is appointed except for one or more of the causes named in the applicable statute.''
\end{quote}
\noindent---Justice Sutherland, \href{https://supreme.justia.com/cases/federal/us/295/602/}{\textit{Humphrey's Executor v. United States}, 295 U.S.\ 602, 631--32 (1935)} (\href{https://www.courtlistener.com/opinion/102475/humphreys-v-united-states/}{CourtListener}).

\end{document}
